\chapter{Descripción del problema}

Este capítulo plantea el problema de investigación, su formulación y las cuatro preguntas de sistematización que guían el desarrollo experimental. También expone el objetivo general y los cuatro objetivos específicos aprobados en el anteproyecto para estructurar las etapas del sistema. Finalmente, presenta la justificación sanitaria, técnica y de ingeniería junto con las delimitaciones del prototipo en nivel de madurez tecnológica (TRL) 4~\cite{amaral2025esmo,vasey2022decideai}.

\section{Formulación}

El melanoma cutáneo registró 337.969 casos nuevos y 58.398 muertes estimadas en el mundo durante 2024~\cite{sung2026globocan,iarc2026melanoma}. En la práctica clínica, la dermatoscopia mejora la precisión diagnóstica de los especialistas y el diagnóstico definitivo exige confirmación histopatológica~\cite{amaral2025esmo}. Aunque se han entrenado clasificadores sobre las 10.015 imágenes de HAM10000~\cite{tschandl2018ham10000}, su evaluación se distorsiona cuando imágenes de una misma lesión quedan en entrenamiento y prueba~\cite{kapoor2023leakage,abhishek2025quality}. Además, gran parte de los desarrollos permanecen como experimentos de laboratorio sin integrarse en aplicativos web accesibles~\cite{farias2026}. A partir de estas necesidades metodológicas y de despliegue, el anteproyecto definió la siguiente pregunta de investigación:

\begin{quote}
¿Cómo desarrollar un sistema de apoyo al diagnóstico de melanoma, basado en redes neuronales convolucionales y técnicas de \textit{transfer learning}, integrado en un aplicativo web que permita la clasificación de imágenes de lesiones cutáneas y obtenga una predicción de riesgo de forma accesible?
\end{quote}

\section{Sistematización}

Para descomponer la pregunta de investigación, el problema se estructuró en cuatro subpreguntas alineadas con las etapas del proyecto. Estas subpreguntas abarcan la gestión de las 10.015 imágenes de HAM10000~\cite{tschandl2018ham10000}, la implementación de arquitecturas con aprendizaje por transferencia~\cite{zhang2024dive}, la selección del clasificador~\cite{maierhein2024metrics} y su despliegue web~\cite{ning2024onnxweb}. A continuación se transcriben textualmente las cuatro subpreguntas formuladas en el anteproyecto:

\begin{enumerate}
  \item ¿Cómo gestionar una base de datos compuesta por imágenes cutáneas que sea etiquetada y que permita la identificación de melanoma?
  \item ¿Cuáles arquitecturas de aprendizaje profundo se pueden implementar para la clasificación de melanoma a partir de imágenes cutáneas?
  \item ¿Cómo seleccionar el mejor modelo de aprendizaje profundo para la identificación de me\-la\-no\-ma?
  \item ¿Cómo integrar el modelo de clasificación seleccionado en un aplicativo web funcional que permita a usuarios cargar imágenes cutáneas y obtener un resultado de identificación?
\end{enumerate}

\section{Objetivos}

Los objetivos del trabajo establecen el propósito general y las cuatro metas operativas correspondientes a las preguntas de sistematización de la Sección~2.2. Su cumplimiento abarca desde el procesamiento del conjunto HAM10000~\cite{tschandl2018ham10000} hasta la verificación del aplicativo web con inferencia local~\cite{ning2024onnxweb,wang2025browser}. A continuación se transcriben de forma textual el objetivo general y los cuatro objetivos específicos del anteproyecto aprobado.

\subsection{Objetivo general}

Desarrollar un sistema de detección de cáncer de piel mediante técnicas de \textit{deep learning}, integrado en un aplicativo web, analizando imágenes de lesiones cutáneas y generando una predicción como apoyo al diagnóstico preventivo.

\subsection{Objetivos específicos}

\begin{enumerate}
  \item Gestionar una base de datos compuesta por imágenes cutáneas etiquetadas aplicando limpieza, preprocesamiento y aumento de datos para el entrenamiento de métodos de aprendizaje profundo.
  \item Implementar arquitecturas basadas en redes neuronales convolucionales (CNN) para la clasificación de melanoma a partir de imágenes cutáneas.
  \item Seleccionar el mejor modelo de aprendizaje profundo para la identificación de melanoma mediante la evaluación de los métodos con métricas estándar de clasificación.
  \item Integrar el modelo seleccionado en un aplicativo web funcional que permita a los usuarios cargar imágenes cutáneas y obtener un resultado de identificación de melanoma.
\end{enumerate}

\section{Justificación}

Desde la perspectiva sanitaria, las 58.398 muertes y los 337.969 casos nuevos estimados en el mundo en 2024 evidencian la carga del melanoma cutáneo~\cite{sung2026globocan,iarc2026melanoma}. En Colombia, la Cuenta de Alto Costo informó que el 67,4\,\% de los casos diagnosticados en el periodo 2024 correspondió a estadios I y II~\cite{cac2025melanoma}. En ese mismo reporte nacional se registraron 11.064 casos atendidos de carácter preliminar, previos a auditoría~\cite{cac2025melanoma}. La finalidad de este prototipo se alinea conceptualmente con el ODS 3 (Salud y bienestar), que plantea fortalecer la prevención y la capacidad de los sistemas sanitarios~\cite{un2024sdg}. Esta relación expresa la orientación del sistema como herramienta de apoyo y no representa una reducción de mortalidad ni un impacto clínico demostrado~\cite{vasey2022decideai}.

En el plano técnico y metodológico, la literatura ha documentado que la fuga de información produce estimaciones optimistas en 17 campos científicos~\cite{kapoor2023leakage}. En dermatología computacional, el conjunto HAM10000 incluye 1.956 lesiones con múltiples fotografías, las cuales suman 4.501 de sus 10.015 imágenes~\cite{tschandl2018ham10000}. Auditorías recientes sobre derivados de HAM10000 demostraron que partir los datos por imagen reparte fotografías de una misma lesión entre entrenamiento y prueba~\cite{abhishek2025quality}. Por ello, este trabajo aplicó una partición por lesión (5.336 lesiones en entrenamiento, 1.067 en validación y 1.067 en prueba) y midió la diferencia frente a partir por imagen. Además, el modelo seleccionado se evaluó sin reajustar sus umbrales sobre 1.512 imágenes de ISIC 2018 sin lesiones compartidas con el entrenamiento~\cite{codella2019isic2018}.

Desde la ingeniería de software, la integración del clasificador mediante ONNX Runtime Web ejecuta el preprocesamiento y la inferencia directamente en el navegador del cliente~\cite{ning2024onnxweb}. Esta arquitectura evita mantener servidores con unidad de procesamiento gráfico (GPU) y procesa cada imagen en 59\,ms en promedio en Chromium sobre unidad central de procesamiento (CPU). De igual forma, las imágenes dermatoscópicas no abandonan el equipo del usuario durante el análisis~\cite{ning2024onnxweb}. Frente al planteamiento inicial del anteproyecto, se precisa que el sistema no está destinado a que la población general evalúe lesiones sin intervención clínica. El aplicativo requiere imágenes adquiridas con dermatoscopio y entrega una predicción de apoyo que no reemplaza el criterio médico ni la confirmación histopatológica~\cite{amaral2025esmo}.

\section{Delimitaciones y alcance}

El sistema desarrollado corresponde a un prototipo validado en entorno de laboratorio (TRL 4) mediante evaluación retrospectiva sobre conjuntos públicos~\cite{tschandl2018ham10000,codella2019isic2018}. Aunque el anteproyecto preveía una ejecución restringida a un entorno local, el aplicativo final se desplegó como sitio web estático en Vercel (\pendiente{URL}). Este cambio de alcance fue posible porque el modelo exportado a ONNX (12,9\,MB) y el preprocesamiento en OpenCV.js se ejecutan íntegramente en el navegador del usuario~\cite{ning2024onnxweb}. De este modo, el despliegue público no requiere infraestructura de cómputo en servidor ni transmisión de imágenes por la red.

El alcance del modelo se circunscribe a imágenes dermatoscópicas clasificadas en las siete categorías pigmentadas definidas en HAM10000~\cite{tschandl2018ham10000}. El clasificador no fue entrenado ni evaluado con fotografías clínicas tomadas con teléfonos móviles ni incluye metadatos del paciente como edad, sexo o localización anatómica. El proyecto no comprende ensayos clínicos prospectivos con pacientes~\cite{marchetti2023proveai,heinlein2024multicenter}, estudios de integración en flujos asistenciales según la guía DECIDE-AI~\cite{vasey2022decideai} ni módulos de interoperabilidad hospitalaria. En consecuencia, toda salida del aplicativo debe interpretarse como información de apoyo sujeta a valoración dermatológica e histopatológica~\cite{amaral2025esmo}.
